% !TeX root = ../初等数学笔记.tex
% 由主文档通过 \input 载入；请编译主文档。
\setcounter{chapter}{23} % 保留原书章号。
\chapter{数学建模、综合方法与后续学习}
\label{chap:em-modeling}
综合应用要求在不同数学表示之间作选择，并把计算结果送回原问题检验。一个模型可以在数学上完全正确，却因假设不合实际而失去解释力；因此变量、约束和适用范围与求解过程同样重要。

\section{从问题到变量、约束与目标}
\emph{数学建模}（mathematical modeling）包括识别对象、提出假设、定义变量、建立关系、求解与检验。未知量是需要确定的值，参数描述外部条件，约束规定允许的取值，目标则说明怎样比较方案。方程描述相等关系，不等式描述容量、范围或预算，整数约束描述不可分割的数量。

例如行程关系 $s=vt$ 依赖速度在所讨论时段恒定；混合物守恒需要所计物质在过程中不被产生或损失。不同单位应先统一，出现负长度、概率超过 $1$ 或应为整数的数量出现小数时，应检查模型与筛选，而非机械接受代数输出。模型假设改变时，原解法可能仍正确，但已经回答另一个问题。

\section{函数、方程与不等式的综合使用}
用长度为 $L>0$ 的围栏围成靠直墙的矩形，墙边无需围栏。令垂直墙边长为 $x$、平行墙边长为 $y$，则 $2x+y=L$、$0<x<L/2$，面积为
\[
 S=x(L-2x)=\frac{L^2}{8}-2\left(x-\frac L4\right)^2.
\]
所以最大面积为 $L^2/8$，在 $x=L/4,y=L/2$ 取得。也可把约束看成 $2x+y=L$，对两正数 $2x,y$ 用均值不等式得到 $2xy\le L^2/4$。两方法分别从顶点与平衡取等解释同一结构。

若 $L=19$ 且两边长须为正整数，连续最优点 $x=19/4$ 不可行。配方说明只需比较允许整数中最接近它的 $4,5$：对应 $(x,y)=(4,11),(5,9)$，面积为 $44,45$，故整数最优方案是 $(5,9)$。不能独立地把连续解各坐标四舍五入，因为那可能破坏 $2x+y=19$。

\section{几何、三角与坐标的综合使用}
考虑水平地面上的竖直目标，忽略观测者高度。在同一直线上靠近与远离目标的两站间距为 $d>0$，仰角分别为 $\alpha,\beta$，满足 $0<\beta<\alpha<\pi/2$。设近站到目标底部距离为 $x>0$、目标高为 $h$，则
\[
 h=x\tan\alpha=(x+d)\tan\beta.
\]
消元得
\[
 x=\frac{d\tan\beta}{\tan\alpha-\tan\beta},\qquad
 h=\frac{d\tan\alpha\tan\beta}{\tan\alpha-\tan\beta}.
\]
若 $d=10,\alpha=45^\circ,\beta=30^\circ$，则 $h=5(\sqrt3+1)$。分母为正与近站仰角较大相一致；若两角非常接近，小测量误差可能显著影响结果，应保留角度误差范围。观测者高度、地面坡度或目标倾斜存在时，须相应修改几何关系。

\section{随机模拟、数据与模型检验}
给定有限概率模型、试验次数 $N$ 和目标事件 $A$，模拟按模型独立生成结果，以指示变量 $I_j$ 记录第 $j$ 次是否发生 $A$，输出经验频率 $\widehat p=N^{-1}\sum_j I_j$。设模型概率为 $p$，则
\[
 \mathbb E[\widehat p]=p,\qquad
 \operatorname{Var}(\widehat p)=\frac{p(1-p)}N.
\]
独立性用于方差计算，伪随机程序也只有在被合理用作该模型的近似生成器时，才支持这种解释。

对任意 $\varepsilon>0$，由于偏差超过 $\varepsilon$ 时平方偏差至少为 $\varepsilon^2$，
\[
 \varepsilon^2\mathbb P(|\widehat p-p|\ge\varepsilon)
 \le\mathbb E[(\widehat p-p)^2]
 =\frac{p(1-p)}N,
\]
故超出误差阈值的概率至多为 $1/(4N\varepsilon^2)$。这是一个可能较松的理论保证，不是某次模拟必然达到的精度。例如 $N=10000,\varepsilon=0.05$ 时上界为 $0.01$。模拟可以发现错误和形成猜测，却不能替代普遍成立的数学证明。

\section{通向高等数学的概念桥梁}
初等数学已经产生了高等数学的主要问题：有限和如何走向无穷和，平均变化如何走向瞬时变化，直线逼近如何刻画曲线，方程组的消元怎样推广到高维。数学分析研究极限、连续、微分与积分，线性代数研究向量与线性映射，概率论把有限模型推广到更一般的随机空间。

以自然对数的底数为例，令 $e_n=(1+1/n)^n$。对 $n$ 个 $1+1/n$ 和一个 $1$ 使用算术--几何平均不等式，可得
\[
 e_n<\left(1+\frac1{n+1}\right)^{n+1}=e_{n+1}.
\]
又二项式定理给出 $e_n\le\sum_{k=0}^n1/k!$。当 $k\ge2$ 时 $k!\ge2\cdot3^{k-2}$，所以对 $n\ge2$，
\[
 e_n\le2+\frac12\sum_{j=0}^{n-2}3^{-j}<\frac{11}4.
\]
故数列递增有界，其极限定义为 $e$，并有 $2<e<3$，从而底为 $e$ 的对数 $\ln x$ 有合法定义。证明 $(e^x)'=e^x$ 等性质还需要建立导数及相应极限定理，本书没有因定义了 $e$ 就提前假定这些结论。这个例子说明，数列、计数、不等式与完备性如何共同支持一个新的数学对象。

\section{综合习题、表达训练与结课自测}
\begin{example}
一个模型给出面积上界 $M$，但所有达到等号的参数都违反原约束。能否宣布最大面积为 $M$？
\end{example}
\begin{solve}
不能。已证的只是每个合法方案的面积不超过 $M$；还需证明某个合法方案达到它，或重新寻找可达的更小上界。若能构造合法方案使面积任意接近 $M$，则 $M$ 是上确界，但仍可能没有最大值。连续模型与整数模型、开区间与闭区间之间的差异常由此产生。
\end{solve}
\begin{enumerate}
\item 靠墙矩形围栏长为 $20$，求最大面积和边长。答：$50$，垂直墙边长 $5$、平行墙边长 $10$。
\item 在上述独立模拟模型中，使误差超过 $0.1$ 的概率上界不大于 $0.05$，次数至少应满足什么条件？答：$N\ge500$；这是充分保证，不一定是最小实际所需次数。
\item 用同一道几何问题比较综合法、三角法与坐标法，逐项说明它们使用的对象、假设与取等或存在条件。评价解答应看论证与适用范围，而非只看计算步数。
\end{enumerate}

